\documentclass[11pt]{article}
\usepackage[letterpaper,margin=1in]{geometry}
\usepackage[T1]{fontenc}
\usepackage{lmodern,microtype,amsmath,amssymb,amsthm,xcolor,tcolorbox}
\tcbuselibrary{skins}
\usepackage[hidelinks]{hyperref}
\definecolor{routegreen}{HTML}{128675}
\definecolor{ink}{HTML}{213337}
\newtcolorbox{takeaway}{enhanced,colback=routegreen!5!white,
 colframe=routegreen!5!white,colupper=ink,boxrule=0pt,
 borderline west={2pt}{0pt}{routegreen},arc=2mm,boxsep=0pt,
 left=4mm,right=4mm,top=2.5mm,bottom=2.5mm,
 before skip=7pt,after skip=7pt,fontupper=\sffamily\small}
\newtheorem{theorem}{Theorem}
\title{Shortest-Path Error from an Edge-Cost Floor}
\author{Jonathan R. Landers}
\date{November 22, 2026}
\makeatletter
\renewcommand{\@maketitle}{\begin{center}
 {\LARGE\@title\par}\vskip .6em
 {\large\@author\par}\vskip .3em
 {\@date\par}\end{center}\vskip .4em}
\makeatother
\begin{document}
\maketitle

A routing model may assign very small costs to some edges. Suppose we
raise every edge below a chosen threshold $k$ up to $k$, then compute a
shortest path. How much can the reported distance increase, and how good
is the selected route when measured with the original costs?

\section*{1. The rounding rule}
Let $G=(V,E)$ be a finite directed or undirected graph with $n\ge2$
vertices and finite costs $w(e)\ge0$. Fix distinct vertices $s,t$ with a
route from $s$ to $t$, and choose $k>0$. Define
\begin{equation}\label{eq:floor}
 \widetilde w(e)=\max\{w(e),k\},\qquad
 \delta(e)=\widetilde w(e)-w(e)=\max\{0,k-w(e)\}.
\end{equation}
Thus $0\le\delta(e)\le k$. For any path $P$, write
\[
 L(P)=\sum_{e\in P}w(e),\qquad
 D(P)=\sum_{e\in P}\delta(e),\qquad
 \widetilde L(P)=L(P)+D(P).
\]
Let $P^*$ be a simple original shortest path with cost $d=L(P^*)$.
Let $\widehat P$ minimize the modified cost, and put
$\widetilde d=\widetilde L(\widehat P)$. Nonnegative original weights
allow a simple $P^*$ by removing cycles.

\section*{2. Two errors, one bound}
\begin{theorem}
If $P^*$ contains $q$ edges with original cost below $k$, then
\begin{align}
 0\le\widetilde d-d&\le D(P^*)\le qk,\label{eq:distance}\\
 0\le L(\widehat P)-d&\le D(P^*)-D(\widehat P)\le D(P^*)\le qk.
 \label{eq:route}
\end{align}
\end{theorem}
\begin{takeaway}
Each rounded edge on an original optimal route adds at most $k$ to the
error bound, covering both the reported distance and the chosen route.
\end{takeaway}
\begin{proof}
Every path becomes at least as expensive, so $\widetilde d\ge d$.
Optimality of $\widehat P$ under the modified costs gives
\[
 L(\widehat P)+D(\widehat P)
 =\widetilde d\le\widetilde L(P^*)=d+D(P^*).
\]
This proves both upper bounds; $L(\widehat P)\ge d$ proves the second
lower bound. Only the $q$ rounded edges contribute to $D(P^*)$, each
by at most $k$.
\end{proof}

\newpage
\section*{3. A bound available before solving the problem}
A simple path has at most $n-1$ edges, so $q\le n-1$ and
\begin{equation}\label{eq:universal}
 \boxed{\quad
 0\le\widetilde d-d\le(n-1)k,\qquad
 0\le L(\widehat P)-d\le(n-1)k.
 \quad}
\end{equation}
For a desired additive tolerance $\varepsilon>0$, choosing
$k\le\varepsilon/(n-1)$ guarantees both errors are at most $\varepsilon$.
The path-specific bound \eqref{eq:distance} is sharper when few edges
are rounded or their costs are close to $k$.

The reported-distance bound is attained by a chain of $n-1$ zero-cost
edges: its original distance is zero and its modified distance is
$(n-1)k$. That chain still selects the same route, so its route error is
zero. Reported-distance inflation and route error measure different effects.

\section*{4. A route can change}
Take three vertices $s,u,t$ with costs $w(s,u)=w(u,t)=0.1$ and
$w(s,t)=1$. The original optimum uses $s\to u\to t$ and has cost $0.2$.
With $k=1$, the two-edge route costs $2$, while the direct edge costs $1$.
The modified optimum therefore selects $s\to t$. Its reported distance
and its original cost are both $1$, giving errors of $0.8$.

Replacing $0.1$ by any $0<\eta<1/2$ gives a selected-route ratio
$L(\widehat P)/d=1/(2\eta)$, which grows arbitrarily large as
$\eta\to0$. The additive bound holds throughout. A relative bound can
use an additional lower scale: if every edge satisfies $w(e)\ge a>0$,
then with $c=\max\{1,k/a\}$,
\[
 \widetilde w(e)\le c\,w(e)
 \quad\Longrightarrow\quad
 L(\widehat P)\le\widetilde d\le c\,d.
\]

\section*{5. What the floor changes}
The rule \eqref{eq:floor} changes edge costs while preserving vertex
locations and graph connectivity. Two locations one metre apart can be
charged ten metres for their connecting edge; their physical separation
remains one metre. Bounds based on geometric packing use actual vertex
separation, whereas the bounds here follow from accumulated cost increases.

The floor also gives a direct bound on hop count: any path with modified
cost at most $U$ contains at most $\lfloor U/k\rfloor$ edges, since each
edge costs at least $k$. This can limit the depth of a search with a known
cost budget; the number of possible routes also depends on branching.

\begin{takeaway}
Raising small costs to $k$ gives an explicit additive error budget and a
hop-count bound for cost-limited paths; the selected route can be evaluated
afterward using the original edge costs.
\end{takeaway}
\end{document}
